\section{线性规划与二次规划的复盘}

\label{sec:8-2}

第 \ref{sec:7-2} 节已经给出了抽象锥规划标准型，第 \ref{sec:7-4} 节则给出了无穷维 KKT 系统。本节要做的，不是重新讲一遍教科书里的 LP/QP 基础，而是把这些抽象对象逐项坍缩回矩阵语言。也正因此，线性规划和二次规划在这里的角色更像“坐标化译本”：同一个原--对偶结构，在有限维里被写成矩阵、向量和互补松弛式。

\subsection{线性规划与其对偶}

\begin{definition}[有限维线性规划]\label{def:linear-program-finite}
设 $c\in \mathbb R^n$，$A\in \mathbb R^{p\times n}$，$G\in \mathbb R^{m\times n}$，$b\in \mathbb R^p$，$h\in \mathbb R^m$。称问题
\[
\min_{x\in \mathbb R^n}\ c^\top x
\qquad \text{subject to}\qquad
Ax=b,\quad Gx\le h
\]
为一个\emph{线性规划}。
\end{definition}

\begin{proposition}[LP 的矩阵对偶形式]\label{prop:lp-dual-matrix-form}
定义~\ref{def:linear-program-finite} 的线性规划，其拉格朗日函数可写为
\[
L(x,\nu,\lambda)=c^\top x+\nu^\top(Ax-b)+\lambda^\top(Gx-h),
\qquad \lambda\ge 0.
\]
相应对偶问题为
\[
\max_{\nu\in \mathbb R^p,\ \lambda\in \mathbb R_+^m}\ -b^\top \nu-h^\top \lambda
\]
subject to
\[
c+A^\top \nu+G^\top \lambda=0.
\]
\end{proposition}

\begin{proof}
这正是第 \ref{sec:7-2} 节定义~\ref{def:abstract-cone-program} 在 $K=\mathbb R_+^m$ 且等式约束单独列出的有限维形式。对偶函数定义为
\[
q(\nu,\lambda):=\inf_{x\in \mathbb R^n}L(x,\nu,\lambda).
\]
若
\[
c+A^\top \nu+G^\top \lambda\neq 0,
\]
则 $L(x,\nu,\lambda)$ 对 $x$ 是非常数仿射函数，其下确界为 $-\infty$。只有当
\[
c+A^\top \nu+G^\top \lambda=0
\]
时，对偶函数才有限，此时
\[
q(\nu,\lambda)=-b^\top \nu-h^\top \lambda.
\]
再结合 $\lambda\ge 0$，便得到所述矩阵对偶形式。
\end{proof}

\begin{theorem}[LP 的矩阵型 KKT 条件]\label{thm:lp-kkt-matrix}
设线性规划存在最优解，并满足有限维 Slater 型严格可行性。则 $x^\star$ 为原问题最优解、$(\nu^\star,\lambda^\star)$ 为对偶最优解，当且仅当
\[
Ax^\star=b,\qquad Gx^\star\le h,
\]
\[
\lambda^\star\ge 0,
\]
\[
c+A^\top \nu^\star+G^\top \lambda^\star=0,
\]
\[
\lambda_i^\star(h_i-g_i^\top x^\star)=0,\qquad i=1,\dots,m.
\]
\end{theorem}

\begin{proof}
这正是第 \ref{sec:7-4} 节定理~\ref{thm:infinite-dimensional-kkt} 在有限维 Euclidean 空间中的坍缩。原可行性与对偶可行性分别对应前两组和第三组条件；驻点条件由
\[
0\in \partial(c^\top x)+A^\top \nu^\star+G^\top \lambda^\star
\]
简化为线性方程
\[
c+A^\top \nu^\star+G^\top \lambda^\star=0;
\]
互补松弛则由
\[
\langle \lambda^\star,h-Gx^\star\rangle=0
\]
在逐项非负锥上的配对写成逐坐标乘积等于零。
\end{proof}

\subsection{二次规划}

\begin{definition}[有限维二次规划]\label{def:quadratic-program-finite}
设 $H\in \mathbb S^n_+$ 为对称半正定矩阵，$c\in \mathbb R^n$。称问题
\[
\min_{x\in \mathbb R^n}\ \frac12 x^\top Hx+c^\top x
\qquad \text{subject to}\qquad
Ax=b,\quad Gx\le h
\]
为一个\emph{凸二次规划}。
\end{definition}

\begin{proposition}[QP 的矩阵型 KKT 系统]\label{prop:qp-kkt-matrix}
在定义~\ref{def:quadratic-program-finite} 中，若满足有限维严格可行性，则 $x^\star$ 为最优解，当且仅当存在 $(\nu^\star,\lambda^\star)$ 使
\[
Ax^\star=b,\qquad Gx^\star\le h,\qquad \lambda^\star\ge 0,
\]
\[
Hx^\star+c+A^\top \nu^\star+G^\top \lambda^\star=0,
\]
\[
\lambda_i^\star(h_i-g_i^\top x^\star)=0,\qquad i=1,\dots,m.
\]
\end{proposition}

\begin{proof}
由 $H\succeq 0$，目标函数是凸且可微，其梯度为
\[
\nabla\!\left(\frac12 x^\top Hx+c^\top x\right)=Hx+c.
\]
因此第 \ref{sec:5-2} 节的一阶微分语言和第 \ref{sec:7-4} 节的 KKT 系统结合后，驻点条件正好化为
\[
Hx^\star+c+A^\top \nu^\star+G^\top \lambda^\star=0.
\]
其余条件与线性规划完全一致，故结论成立。
\end{proof}

\begin{remark}
命题~\ref{prop:qp-kkt-matrix} 说明，第 \ref{sec:7-4} 节中的抽象乘子系统在有限维里会同时吸收第 \ref{sec:5-2} 节的梯度与第 \ref{sec:5-3} 节的次微分语言。LP 对应的是线性目标的次梯度为常向量，QP 对应的是二次目标的梯度为仿射映射；而结构本身并没有变化。
\end{remark}

\subsection{典型例子}

\begin{example}[多面体上线性目标]
线性规划中的可行集是多面体，目标则是线性泛函。把这个例子放在这里，是为了把第 \ref{sec:4-4} 节极点与支撑超平面的语言回收到最经典的 LP 场景：抽象分离与极点理论，在有限维里正好对应多面体几何。
\end{example}

\begin{example}[标准二次规划]
当 $H$ 为对称半正定矩阵时，定义~\ref{def:quadratic-program-finite} 给出的 QP 同时体现了第 \ref{sec:5-2} 节的梯度、第 \ref{sec:7-2} 节的标准型和第 \ref{sec:7-4} 节的 KKT 条件。把这个例子放在这里，是为了让读者看到 Hessian、乘子和互补松弛在有限维中如何自然拼接起来。
\end{example}

\begin{example}[均值--方差投资组合]
均值--方差组合优化可写为
\[
\min_x \frac12 x^\top \Sigma x-\mu^\top x
\qquad \text{subject to}\qquad
\mathbf 1^\top x=1,\quad x\ge 0.
\]
把这个例子放在这里，是为了给第 \ref{chap:11} 章连续时间模型一个静态有限维基线，同时说明 QP 的矩阵型 KKT 公式具有直接的经济解释。
\end{example}

\subsection*{有限维对照}

Boyd 书中 LP 和 QP 的拉格朗日对偶与 KKT 条件通常直接写成矩阵形式，这会给人一种错觉：似乎这些结论只属于线性代数世界。本节的作用正相反。通过回链定义~\ref{def:abstract-cone-program}、定理~\ref{thm:infinite-dimensional-kkt} 与推论~\ref{cor:kkt-via-fenchel-subdifferential}，我们看到矩阵公式只是抽象 KKT 在 $\mathbb R^n$ 中的坐标表示。

\subsection*{习题（待补）}

\begin{exercise}[LP 对偶占位]
补一题：从定义~\ref{def:linear-program-finite} 出发，直接计算线性规划的对偶函数并推出命题~\ref{prop:lp-dual-matrix-form}。
\end{exercise}

\begin{exercise}[QP KKT 桥接题占位]
补一题：对一个带不等式约束的凸二次规划写出命题~\ref{prop:qp-kkt-matrix} 的全部条件，并说明哪些部分来自第 \ref{sec:7-4} 节的抽象 KKT 系统。
\end{exercise}
